\chapter{Blockchain Beyond Finance: Real-World Applications}
\label{chap:real-world-blockchain}

The previous chapter framed Ethereum and AAVE through the lens of
decentralised finance: how users acquire tokens, what assets exist, why
overcollateralised borrowing is useful, and how DeFi yield compares with a
traditional bank account. That framing is incomplete. The same underlying
infrastructure -- public smart contracts, deterministic execution,
transferable on-chain claims -- has been used for purposes that have nothing
to do with margin loans or yield farming. This chapter looks at two of those
threads. The first is the tokenisation of \emph{real-world assets} (RWAs),
where Treasury bills and apartments are wrapped into transferable tokens.
The second is the use of blockchain rails in low-income, low-trust
environments, illustrated by a concrete case study from Kenyan agriculture.

The aim is not to argue that blockchain is the right answer everywhere. It
is to show two settings where the technology is already deployed at
measurable scale, and to be honest about what it does and does not solve in
each case.

\section{Tokenised Real-World Assets}

A real-world asset, in the on-chain sense, is anything whose ownership lives
primarily off-chain -- a Treasury bill, a corporate bond, a piece of real
estate, a barrel of oil -- but whose right of ownership is represented by a
token on a public blockchain. The token does not replace the underlying
legal claim; it is a programmable wrapper around it, issued by a regulated
entity that holds the off-chain asset in custody. The benefit the token
offers is the same set of properties that ETH itself offers: it can be
transferred 24/7, it can be used as collateral in a smart contract, it can
be split across an arbitrary number of investors, and its movements are
publicly auditable.

\subsection{Tokenised government bonds}

The first RWA category to find clear product-market fit has been
short-dated government debt. A token issuer holds Treasury bills in
custody off-chain and, in exchange, mints an ERC-20 token whose price is
pegged to one dollar and which pays the underlying yield directly to its
holder. Once on-chain, the token behaves like any other ERC-20: it can be
sent in a single transaction at any time of day, used as collateral in
another smart contract, or split across many holders without any paperwork.
BlackRock's BUIDL fund, launched on Ethereum in March 2024, is the
best-known instance and reached several hundred million dollars in
deposits within weeks of launch \cite{blackrockBuidl2024,coindeskBuidl2024}.

The technically interesting property is composability. Once the bond is
represented as a token, the rest of the on-chain ecosystem can use it as
an input without any further integration work. A smart account, an AMM, or
an AAVE-style lending pool can hold or accept the token by virtue of the
ERC-20 interface alone -- the same property this thesis relies on for
every interaction with the AAVE pool. A paper-based claim to a Treasury
bill cannot be plugged into any of those systems.

\subsection{Tokenised real estate}

A more ambitious RWA category is residential real estate. The pattern is
the same as for bonds: an off-chain entity holds the legal title to a
property and issues ERC-20-style tokens that represent fractional
ownership of it; rental income is distributed to token holders on-chain.
Tokens are typically priced at a few tens of dollars, so the barrier to
real-estate investment is reduced by several orders of magnitude.

The interesting twist with real estate is that it exposes the central
limitation of any RWA design. A 2025 investigation into RealT, one of the
larger U.S.\ tokenisation platforms, found that the company had issued
tokens for properties whose deeds were never transferred to the issuing
entity and that some underlying buildings were uninhabitable
\cite{realtInvestigation2025}. The smart contract was working exactly as
specified: tokens were minted, rental payments were distributed. What it
could not verify was whether the off-chain reality matched the on-chain
representation.

\subsection{What this means for the rest of the thesis}

The pattern across both examples is the same. Tokenisation makes a claim
liquid, divisible, and programmable. It does not, on its own, improve the
quality of the underlying claim. Where the off-chain side is mechanical
and well-regulated -- a Treasury bill held by an established custodian --
the wrapper adds genuine value, because every property of the asset is
already verifiable by parties who do not need to trust each other. Where
the off-chain side depends on a single operator's honesty -- a startup
buying houses, a private credit fund, a commodities trader -- the wrapper
inherits whatever weakness exists in that operator.

This is not a flaw in the technology. It is the consequence of a clean
boundary: a smart contract can only enforce the properties of data that
lives on-chain. Whether the off-chain world has supplied that data
faithfully is a question of legal structure, audit, and institutional
trust, not of cryptography. The same observation will reappear in the
agricultural case studies in the next section, where the off-chain inputs
are weather measurements rather than property deeds, and where the
trustworthiness of those inputs is what determines whether the contract is
useful.

\section{Blockchain in Low-Income, Low-Trust Environments}

A different and arguably more interesting use of the same infrastructure
has emerged in settings where the alternative is not a Bloomberg terminal
but a paper ledger, a verbal agreement, or no service at all. Smallholder
agriculture in sub-Saharan Africa is one of those settings. Roughly 7.5
million smallholder farmers work the land in Kenya alone
\cite{etheriscLcc2023}; most are excluded from formal insurance, from
traceable market data, and often from any record of the transactions they
make with intermediaries. Two projects -- both running on public Ethereum
or EVM-compatible chains, both small but well-documented -- give a
concrete sense of what blockchain rails actually do for these users.

\subsection{Parametric crop insurance: ACRE Africa and Etherisc in Kenya}

Crop insurance for smallholders has been a long-running development
challenge. The traditional model -- a claims adjuster visits the farm,
assesses the damage, files paperwork, and triggers a payout weeks later --
does not scale economically when premiums are a few dollars and farms are
distributed across thousands of villages. The result is that farmers who
most need protection from drought or flood cannot access it at a price
that makes sense for either party.

The pilot run by ACRE Africa with Etherisc and Mercy Corps Ventures
replaces the human adjuster with a parametric design. The product, called
Bima Pima, is a weather-index insurance: it pays out when an objective
measurement of weather -- rainfall over the growing season, measured by
satellite -- falls below a defined threshold for the farmer's region. There
is no inspection, no claim form, and no negotiation. The Etherisc Generic
Insurance Framework, deployed as a set of Ethereum smart contracts,
encodes the policy conditions; Chainlink oracles feed in the satellite
rainfall data from the Africa Rainfall Climatology dataset; and when the
contract fires, a payout is delivered to the farmer's M-Pesa mobile money
account \cite{mercyCorpsEtherisc2022}.

The cost structure is what makes the design viable. Farmers pay
approximately fifty cents to one dollar in premiums per bag of seed, which
is bundled into the seed price. A scratch card distributed with the seed
bag carries a code that the farmer texts in by USSD from a feature phone,
which is enough to activate the policy. The whole user-facing flow
requires neither a smartphone nor an internet connection -- which is
critical, since neither is universal in rural Kenya. The blockchain layer
is invisible to the farmer and exists primarily to make the back-end
verifiable: payouts cannot be lost in administrative bottlenecks, and the
relationship between rainfall data and payout decisions is provable
in audit.

The pilot covered roughly 12{,}500 maize, sorghum, and green-gram farmers
in Kenya's 2021 long rains season \cite{mercyCorpsEtherisc2022}. By 2022 a
follow-up cohort of 7{,}000 more farmers had been added through the
Lemonade Crypto Climate Coalition, on the Avalanche chain rather than on
Ethereum, and Etherisc reported having cumulatively insured around
25{,}000 farmers across the ACRE Africa partnership
\cite{etheriscLcc2023}. The numbers are modest against ACRE Africa's
overall reach of 1.7 million smallholders across East Africa, but they are
not pilot-scale either: the product is in repeat use across multiple
growing seasons and has been replicated to additional crops and
geographies.

What does the blockchain layer actually contribute? Removing it would not
make the parametric design impossible: a centralised database run by an
insurer could in principle issue the same payouts. The argument made by
the project teams is that for a population whose primary objection to
formal insurance is mistrust, a system in which the payout rule is
publicly inspectable and the oracle data is signed by a known source is
qualitatively different from a system in which an opaque corporate process
decides whether a claim is honoured. A farmer who has been historically
underpaid by intermediaries does not need to take the protocol's word for
it; the conditions are written in a contract that nobody can edit after
the fact.

\subsection{Supply chain transparency: AgUnity in Kakamega}

A second case study, also from western Kenya, addresses a different
problem: information asymmetry along the agricultural value chain. The
Australian startup AgUnity, working with Virginia Tech's Center for
International Research, Education, and Development and Kenya's Egerton
University, deployed a smartphone application called the AgUnity V3 Super
App in Kakamega County. The app is targeted at smallholder farmers of
African Indigenous Vegetables and at the traders and retailers who buy
from them, and its central feature is a blockchain-backed record of every
transaction in the supply chain
\cite{agunityAIV2022,disruptionAgUnity2022}.

The problem the app addresses is that smallholder farmers historically
operated as price-takers. Brokers would arrive at the farm, name a price,
weigh the produce on their own scales, and pay in cash with no receipt.
Farmers had no record of what they had sold, no information about what
prices were being paid elsewhere in the country, and no leverage in the
next negotiation. The result was systematic underpayment, a high rate of
exploitative brokering, and a strong incentive for farmers to under-invest
in producing quality grades that they could not prove or charge for.

The intervention was a smartphone -- pre-loaded with the AgUnity app and
sometimes co-financed -- given to each actor in the value chain. The app
records each transaction (quantity, grade, price, parties) and writes it
to a blockchain ledger that the farmer, the trader, and other value-chain
participants can all see. The transaction record cannot be retroactively
altered by any single party. Farmers gain a permanent history of their
sales, traders gain a reliable record of the supply they have purchased,
and consumers downstream can verify whether a vegetable graded as
``pesticide-free Grade A'' has actually been recorded as such throughout
its journey \cite{disruptionAgUnity2022}.

The pilot in Kakamega deployed roughly fifty to sixty smartphones to
producers, traders, and retailers over six months in 2021. The reported
outcomes are concrete: smallholder incomes rose because farmers could now
negotiate prices with reference to recent transaction data; weights and
quantities became standardised on kilograms rather than ad-hoc volumes;
post-harvest losses dropped because supply and demand could be matched in
near real time; and reliance on brokers as the only source of market
information decreased \cite{agunityAIV2022,disruptionAgUnity2022}. The
project has since extended a token, AgriUT, used as a reward unit, to
several hundred thousand farmers in the broader region
\cite{disruptionAgUnity2022}.

As with the Etherisc pilot, the blockchain layer is largely invisible to
the user. What the user experiences is a smartphone app that remembers
their transactions and lets them compare prices. What the blockchain
provides is a guarantee that the record they are looking at has not been
silently rewritten by the trader on the other side of the transaction.
That guarantee turns out to be the structurally important property: it is
what makes the difference between yet another corporate database (which
the farmer has no reason to trust more than the broker who is already
cheating them) and a shared ledger whose contents are demonstrably the
same for everyone with a phone.

\subsection{What these case studies show}

Both the parametric insurance pilot and the AgUnity deployment make a
similar argument. The technical contribution of blockchain in these
settings is not throughput, programmability, or yield. It is that the
system's record of what happened is verifiable by parties who do not trust
each other and who would have no other shared point of reference. The
smallholder farmer cannot audit the books of a multinational reinsurer or
of a vegetable broker. They can, however, look at a transaction confirmed
on a public chain.

It is worth being careful with the framing. Neither project replaces the
local institutions on the ground -- the village-based agents who train
farmers on the insurance product, the seed retailers who hand out the
scratch cards, the field officers who explain how to grade a vegetable.
What blockchain enables is a different distribution of trust within those
existing institutions: the parts of the system that handle data and
disbursement become inspectable rather than opaque, and the human-facing
parts get to focus on the work that genuinely requires a human.

The numbers also matter. Twenty-five thousand insured farmers and a few
hundred thousand AgriUT recipients are not, on a continental scale, a
revolution. But they are also not zero, they are not theoretical, and the
products have survived multiple growing seasons -- which is the relevant
bar for an agricultural intervention.

\section{Synthesis}

Tokenised Treasuries and tokenised real estate, smart-contract crop
insurance and on-chain supply-chain records, look superficially very
different. They share a single underlying primitive: a record of an
agreement, written into a smart contract, that all parties can inspect and
that no single party can rewrite. The cases where this primitive delivers
value are the cases where the alternative is a record that one party owns
and the others have to take on faith -- whether that is a sovereign-bond
custody chain, a rental ledger, a claims adjuster's report, or a broker's
verbal price quote.

The cases where blockchain has not yet delivered value are equally
informative. Where the underlying asset is illiquid, opaque, or
operationally complex, tokenisation does not fix it; the RealT
investigation is the cleanest demonstration of that limit. Where a single
trusted intermediary is already cheap and reliable, replacing it with a
smart contract usually makes things slower and more expensive; this is why
no one runs a domestic credit-card network on Ethereum.

The AAVE executor module developed in the rest of this thesis sits firmly
inside the first category. AAVE is itself an on-chain primitive --
positions, debts, and liquidations are public state -- and the module
exposes that primitive to a smart account in a way that other on-chain
components (validators, hooks, paymasters) can plug into. The broader
point of this chapter is that the same trust-redistribution that makes a
parametric insurance payout enforceable for a smallholder farmer in
Kakamega is the trust-redistribution that makes the executor module
useful: in both cases, the contract is the spec, the data is the audit
trail, and no party needs to take any other party's word for what
happened.
